\section{연구 결과}

% --------------------------------------------------------------------
% 이 파일의 역할
% --------------------------------------------------------------------
% 4-Results.tex에는 연구 문제에 대한 답이 되는 연구 결과를 씁니다.
% 표, 그림, 그래프, 통계량, 질적 분석 결과 등을 제시하고,
% 각 결과가 어떤 연구 문제에 대응하는지 분명하게 설명합니다.
% 결과 장은 새로운 주장보다 자료에서 확인된 사실을 중심으로 쓰며,
% 해석과 일반화는 결론 장에서 더 넓게 다룹니다.
%
% --------------------------------------------------------------------
% 연구 결과 장 작성 안내
% --------------------------------------------------------------------
% 이 장에서는 연구 방법에서 수집하고 분석한 자료의 결과를 제시합니다.
% 가장 중요한 원칙은 연구 결과가 연구 문제에 대한 답으로 구성되어야 한다는 점입니다.
% 따라서 결과 장의 하위 절은 가능하면 "연구 문제 1", "연구 문제 2"처럼
% 연구 문제의 순서와 직접 대응하도록 배열하는 것이 좋습니다.
% 각 절에서는 해당 연구 문제에 대해 어떤 답을 얻었는지 먼저 밝히고,
% 그 답을 뒷받침하는 표, 그림, 통계값, 질적 사례를 제시합니다.
% 해석과 논의는 다음 장(논의 또는 결론)에서 자세히 다루고, 여기서는
% 연구 문제에 직접 대응하는 결과를 객관적으로 정리하는 데 초점을 둡니다.
%
% 작성 순서 예시
% 1. 연구 문제 또는 분석 단계별로 하위 절을 나눕니다.
% 2. 각 절의 첫 문단에서 어떤 결과를 제시하는지 간단히 안내합니다.
% 3. 핵심 수치는 표로 정리하고, 변화 추세나 집단 간 차이는 그래프로 제시합니다.
% 4. 표와 그림을 본문에서 반드시 언급합니다.
%    예: 표~\ref{tab:result-example}에 제시한 바와 같이 ...
%    예: 그림~\ref{fig:graph-example}은 집단별 평균 점수의 변화를 보여준다.
% 5. 표나 그림 아래 내용을 그대로 반복하기보다, 독자가 주목해야 할 패턴을 문장으로 설명합니다.
%
% 표 작성 팁
% - 표 제목은 표 위에 둡니다.
% - 단위가 있으면 열 제목이나 표 주석에 표시합니다.
% - 소수점 자릿수는 같은 종류의 값끼리 통일합니다.
% - 너무 긴 표는 부록으로 보내고, 본문에는 핵심 결과만 제시합니다.
% - 표 번호는 직접 입력하지 않습니다. \caption과 \label을 사용하면 LaTeX이 자동으로 번호를 붙입니다.
% - 표를 본문에서 언급할 때는 표~\ref{tab:label-name}처럼 참조합니다.
% - 표 전체 너비는 \linewidth를 기준으로 맞추고, 각 열의 너비는 내용의 길이에 따라 조절합니다.
% - 표 안에는 \ThesisTableStyle을 넣어 글자 크기, 열 간격, 행 간격을 일관되게 맞춥니다.
% - 긴 설명이 들어가는 열은 p{0.25\linewidth}처럼 폭을 지정하고, 숫자 열은 좁게 잡습니다.
%
% 그래프 작성 팁
% - 그림 제목은 그림 아래에 둡니다.
% - 축 이름과 단위를 반드시 표시합니다.
% - 그래프 안의 색상에만 의존하지 말고, 선 모양이나 마커도 구분합니다.
% - 본문에서는 그래프가 보여 주는 핵심 경향을 설명합니다.

연구 결과 장은 과학교육 연구 문제에 대한 답을 자료와 분석 결과로 제시하는 곳이다.
각 절은 연구 문제의 순서와 대응되게 구성하고, 표나 그림을 제시한 뒤에는 학습자의 개념 변화, 탐구 수행 양상, 응답 범주, 집단 간 차이, 수업 장면의 특징처럼 독자가 주목해야 할 패턴과 의미를 본문에서 설명한다.
결과 장에서는 새로운 주장이나 과도한 일반화보다 검사, 설문, 면담, 관찰, 산출물 분석에서 확인된 사실을 객관적으로 정리하는 데 초점을 둔다.
표·그림·수식을 LaTeX으로 작성하는 구체적인 방법(캡션·라벨·참조 명령, 표 서식, 컬러/흑백 관리, TikZ 등)은 \href{../manual/8-KNUE_THESIS_TEMPLATE_USAGE.md}{manual/7 문서}에 정리되어 있으므로 여기서는 반복하지 않는다.

\subsection{연구 문제 1의 결과}

% 이 절에는 첫 번째 연구 문제에 대한 결과를 씁니다.
% 예: 기술통계, 빈도 분석, 사전-사후 검사 결과, 문항별 응답 분포 등
% 연구 문제와 직접 관련된 결과를 먼저 제시합니다.

표~\ref{tab:result-example}은 집단별 사전 검사와 사후 검사 평균 점수를 정리한 예시이다.
실제 논문을 작성할 때는 아래 표의 집단명, 평균, 표준편차, 사례 수를 연구 자료에 맞게 수정한다.

\begin{table}[htbp]
    \centering
    \caption{집단별 사전-사후 검사 점수 예시}
    \label{tab:result-example}
    \ThesisTableStyle
    \begin{tabular*}{\linewidth}{@{\extracolsep{\fill}}p{0.21\linewidth}rrrr@{}}
        \toprule
        집단 & 사례 수 & 사전 평균 & 사후 평균 & 평균 변화 \\
        \midrule
        실험 집단 & 30 & 68.40 & 82.15 & 13.75 \\
        비교 집단 & 28 & 67.85 & 73.20 & 5.35 \\
        \bottomrule
    \end{tabular*}
\end{table}

표~\ref{tab:result-example}에서 실험 집단은 사전 평균 68.40점에서 사후 평균 82.15점으로 증가하였다.
비교 집단도 점수가 증가하였으나, 평균 변화 폭은 실험 집단에서 더 크게 나타났다.
다만 이 차이가 통계적으로 유의한지 판단하려면 $t$ 검정이나 분산분석 같은 추리통계 결과를 함께 제시해야 한다.

\subsection{연구 문제 2의 결과}

% 이 절에는 두 번째 연구 문제에 대한 결과를 씁니다.
% 시간에 따른 변화, 집단 간 비교, 영역별 차이처럼 패턴을 보여 주고 싶을 때는
% 표와 함께 그래프를 사용하면 좋습니다.

그림~\ref{fig:graph-example}은 집단별 평균 점수 변화를 선 그래프로 나타낸 예시이다.
그래프를 사용할 때는 축 이름, 범례, 단위를 명확히 표시한다.

\begin{figure}[htbp]
    \centering
    \begin{tikzpicture}
        \begin{axis}[
            width=0.78\linewidth,
            height=0.42\linewidth,
            font=\footnotesize,
            xlabel={검사 시점},
            ylabel={평균 점수},
            symbolic x coords={사전, 사후},
            xtick=data,
            ymin=60,
            ymax=90,
            legend pos=north west,
            grid=major
        ]
        \addplot+[mark=o] coordinates {(사전,68.40) (사후,82.15)};
        \addlegendentry{실험 집단}
        \addplot+[mark=square] coordinates {(사전,67.85) (사후,73.20)};
        \addlegendentry{비교 집단}
        \end{axis}
    \end{tikzpicture}
    \caption{집단별 평균 점수 변화 예시}
    \label{fig:graph-example}
\end{figure}

% 그래프를 제시한 뒤에는 그래프가 의미하는 바를 간단히 설명합니다.
% 결과 장에서는 "왜 이런 결과가 나왔는지"를 길게 해석하기보다,
% 어떤 차이와 경향이 관찰되었는지를 정확하게 쓰는 것이 좋습니다.

그림~\ref{fig:graph-example}에서 두 집단 모두 사후 점수가 상승하였지만,
실험 집단의 상승 폭이 비교 집단보다 크게 나타나는 경향을 보였다.

두 집단의 평균 변화 폭 차이는 다음과 같은 평균 산출식을 바탕으로 계산하였다.

\begin{equation}
    \bar{x} = \frac{1}{n}\sum_{i=1}^{n} x_i
    \label{eq:mean-example}
\end{equation}

식~\ref{eq:mean-example}에서 \(\bar{x}\)는 표본 평균, \(n\)은 사례 수, \(x_i\)는 \(i\)번째 관측값을 의미한다.

\subsection{추가 분석 결과}

% 필요한 경우 추가 분석 결과를 제시합니다.
% 예: 하위 영역별 결과, 문항별 결과, 질적 응답 범주, 면담 분석 결과,
% 상관분석 또는 회귀분석 결과 등
%
% 질적 결과를 제시할 때는 범주명, 의미, 대표 인용문을 함께 제시하면 좋습니다.
% 단, 참여자를 식별할 수 있는 이름, 학교명, 연락처 등은 익명화해야 합니다.

결과 장의 마지막에는 연구 문제별 답이 충분히 제시되었는지 확인한다.
다음 장에서는 이 결과를 다시 과학 교수·학습 맥락과 연구 목적에 연결하여 결론, 의의, 제한점, 제언으로 일반화한다.
